\section{Methodology}\label{sec:methodology}

\subsection{Dataset}\label{sec:method-dataset}

The dataset consists of seven non-trivial open-source applications whose maintainers ship their own Kubernetes manifests (Table~\ref{tab:dataset}). A project is eligible if (i) its container images are public and pullable without credentials; (ii) it runs on a local \texttt{kind} cluster \citep{tool_kind} without cloud-specific dependencies; (iii) its manifests may be redistributed (otherwise only the commit is stored); and (iv) it is a real application rather than a single-container toy, with at least two cooperating components, persistent state, or genuine functionality such as a workflow engine or a storefront. The seven span four classes and both manifest intents, so that the production stratum of RQ1 is not empty. A candidate was selected only if its human manifest could be deployed and measured on the evaluation cluster, so that every source of configurations could be scored on all six dimensions. This criterion makes the human runtime outcomes partly true by construction, and they are therefore reported as a reference, not as a test of human against LLM deployability (Section~\ref{sec:res-rq1}).

\begin{table*}[t]
\centering\small
\caption{Evaluated projects. Intent is taken from the repository's own wording about its manifests (evidence recorded per project); size and workloads refer to the shipped human manifest; stars are those of the repository holding the manifest, at the fetch date (October 2026).}\label{tab:dataset}
\begin{tabular}{llllrrrr}
\toprule
Project & Class & Intent & Commit & Components & Lines & Objects & Stars \\ \midrule
n8n (n8n-hosting) & stateful & demo & 7a22e91 & 2 & 250 & 9 & 1,750 \\
Grafana & stateful & production & 3920186 & 1 & 79 & 3 & 77,345 \\
Example Voting App & multi-service & demo & 63e9150 & 5 & 182 & 9 & 5,832 \\
Emojivoto & multi-service & demo & 9799218 & 4 & 210 & 11 & 218 \\
Online Boutique & microservice system & demo & 38e7348 & 12 & 912 & 35 & 21,061 \\
Sock Shop & microservice system & demo & 9dff06f & 13 & 844 & 29 & 3,781 \\
Kubernetes Dashboard & cluster add-on & production & 42deb6b & 2 & 242 & 14 & 15,407 \\
\bottomrule
\end{tabular}
\end{table*}

For each project, a \texttt{docs\_bundle/} is assembled by copying the project's own README, Dockerfile(s) and Docker Compose file verbatim at the pinned commit, with only Kubernetes-specific content redacted: any embedded manifest fragment, any ``run on Kubernetes'' deployment section, and any individual sentence stating a fact that is otherwise only recorded in the human YAML (e.g., a replica count or a NodePort value). Every redaction is logged with its file, line range and reason. No sentence in \texttt{docs\_bundle/} is written by the research team or by any LLM. An internal \texttt{fact\_sheet.md}, never shown to a model, records the application's purpose, each component's image and required environment variables, and justifies the per-project requirement checks (D4) and the functional smoke test (D5), both written during dataset construction; every later change to them is one of the logged corrections of Section~\ref{sec:threats}. Each project's bundle and fact sheet are reviewed by a second author against a fixed checklist before generation.

\subsection{Fairness of the human comparison}\label{sec:method-fairness}

The human developers and the LLMs do not work under identical conditions: the human's goal was often a demo rather than production hardening, the human iterated with real feedback while each LLM makes a single attempt, and the human had the full repository while each LLM sees only \texttt{docs\_bundle/}. The design mitigates this asymmetry in four ways. First, the human manifest is treated as a real-world reference point, and claims are phrased as comparisons with what the project ships, not with human engineers in general. Second, the human comparison (RQ1) uses the \emph{neutral} prompt, which, like the human's situation, carries no explicit push toward hardening; the \emph{production} prompt is analysed separately (RQ5). Third, every human manifest is labelled production- or demonstration-oriented from the repository's own wording, and RQ1 is reported overall and per label. Fourth, no specification is compiled by the research team: an earlier design used a team-written specification that an LLM had helped to compile, which could have favoured that model family, and the final design removes it so that every model receives only the maintainers' own words. The LLM-versus-LLM comparisons (RQ3, RQ5) use identical inputs for every model.

\subsection{Generation}\label{sec:method-generation}

Every model receives the same fixed system prompt, ``You are an experienced DevOps engineer. Respond with Kubernetes YAML manifests only,'' followed by a user prompt that concatenates the files in \texttt{docs\_bundle/} (each preceded by its filename, in a fixed order) and ends in one of two instruction sentences. The \emph{neutral} variant asks to ``write Kubernetes manifests to deploy this application''; the \emph{production} variant asks to ``write production-ready Kubernetes manifests.'' The variants differ in no other respect. Three samples are drawn per (model, project, prompt variant), each from a fresh context with no memory of any other generation. Each generation is stored with its request parameters, the raw response, the extracted YAML (every fenced YAML block, joined in order, or the whole response if there is none) and the extraction status. A reply is scored as returned: one whose YAML block is opened but never closed falls back to the whole response and therefore fails D1, as it would for a user who pasted it.

Table~\ref{tab:models} lists the five models and how each was accessed. For budget reasons, only Gemini and Qwen were accessed through an API (free tiers), and only for them could temperature (0.2) and the output limit (8,192 tokens) be set; the other routes do not expose sampling parameters. Each model is therefore evaluated together with its access route, and the comparisons of RQ3 are comparisons of these model-and-access configurations. Claude was generated by isolated Claude Code subagents rather than a bare API call; the transcripts of all 45 recorded generation sessions (the 42 retained samples and three pilot samples generated under an earlier protocol, which were discarded and regenerated) show that each subagent read only its own prompt file and wrote only its output manifest, and none accessed the human manifests, other models' outputs or the web. The consequences of these differences are discussed in Section~\ref{sec:threats}.

\begin{table*}[t]
\centering\small
\caption{Models and access routes. Temperature and output limit could be set only for the API-generated models. ``Third party'': prompts were sent and replies copied by an operator outside the research team, under a written protocol (Section~\ref{sec:threats}).}\label{tab:models}
\resizebox{\textwidth}{!}{%
\begin{tabular}{llllllll}
\toprule
Model & Identifier as recorded & Tier and reasoning & Access route & Generated (2026) & Temp. & Max tokens & Operator \\ \midrule
Claude & claude-sonnet-5-5 & mid tier; harness default & Claude Code subagent, fresh context & 9 Nov & -- & -- & authors \\
DeepSeek & DeepSeek web chat & chat model; reasoning off & chat.deepseek.com, new chat per prompt & 7--9 Nov & -- & -- & authors \\
Gemini & gemini-3.6-flash & flash tier; hidden reasoning & Gemini API (free tier) & 4--11 Nov & 0.2 & 8,192 & authors \\
GPT & ChatGPT default model & default; medium reasoning & chatgpt.com, temporary chat per prompt & 8--13 Nov & -- & -- & third party \\
Qwen & qwen3.8-27b & 27B open weights; provider default & Groq API (free tier) & 4--11 Nov & 0.2 & 8,192 & authors \\
\bottomrule
\end{tabular}}
\end{table*}

\subsection{K8sGenBench: a six-dimension benchmark}\label{sec:method-benchmark}

Every configuration is evaluated by the same automated pipeline, using deterministic tools (no LLM-as-judge) and, for D1--D3, at least two tools per dimension.

\textbf{D1 -- Correctness.} The manifest is parsed with PyYAML, rejecting duplicate mapping keys as YAML~1.2 and \texttt{kubectl} do, validated against pinned Kubernetes schemas with \texttt{kubeconform} \citep{tool_kubeconform}, and submitted object by object to the API server with \texttt{kubectl apply --dry-run=server}, so that one object the cluster cannot resolve does not mask the validity of the others. Metrics are the parse rate, the schema-pass rate and the share of objects the API server accepts.

\textbf{D2 -- Security.} Configurations are scanned with Kubescape \citep{tool_kubescape} against the NSA-CISA \citep{nsacisa2022k8shardening} and MITRE frameworks and with Trivy's configuration scanner \citep{tool_trivy}, and are dry-run into namespaces that enforce the \texttt{baseline} and \texttt{restricted} Pod Security Standards \citep{k8s_pod_security_standards}. A secret check counts literal values under credential-named keys and known credential token formats, ignoring placeholder values. Findings are also classified into the eleven misconfiguration categories of Rahman et al.\ \citep{rahman2023k8smisconfig}, using the rules of their SLI-KUBE tool applied to every pod and container specification.

\textbf{D3 -- Static reliability.} Configurations are scanned with kube-score \citep{tool_kubescore}, Polaris \citep{tool_polaris} and KubeLinter \citep{tool_kubelinter}, including the six KubeLinter smells used by Zhang et al.\ \citep{zhang2024genaismells}, reported as density per thousand lines. A ten-item Reliability Controls Checklist is applied to every workload: at least two replicas for Deployments, a readiness probe, a liveness probe, a CPU request, a memory request, a memory limit, an explicit rolling-update strategy, a PodDisruptionBudget, pod anti-affinity or a topology-spread constraint, and a pinned image tag. Items that do not apply to a workload kind are skipped, and checklist coverage is the share of applicable items met.

\textbf{D4 -- Requirement satisfaction.} Each project's \texttt{requirements.py} checks facts that are traceable to a file in \texttt{docs\_bundle/}, for example that every documented component is deployed or that a documented required environment variable is set. Facts that are true of the project but cannot be recovered from the bundle (for example, an image registry path that appears only in unbundled build scripts) are reported as \emph{oracle} checks, a diagnostic of what a model knew beyond its input, and excluded from the score. The metric is the share of scored requirements met.

\textbf{D5 -- Deployment.} Each configuration is deployed into a fresh namespace on a local \texttt{kind} cluster and exercised with a project-specific smoke test: a sequence of HTTP requests (and, where the documentation states a data-path effect, a command run inside a data-store container) issued from a client pod inside the cluster to the Service that fronts the documented entry component, each expecting a documented status code and response content. A configuration counts as \emph{fully Ready} only once every workload has held its desired number of Ready replicas for 20~seconds without interruption; the deadline is 300 seconds. Client pods comply with the \texttt{restricted} Pod Security Standard and are placed where the configuration's own NetworkPolicies admit a client, so that a configuration that deliberately restricts ingress is not mistaken for one that fails to serve. When a configuration does not become Ready, the cluster's record is captured at the deadline: per container, the restart count, the last exit code and reason and the last log lines, and per object, the warning events.

\textbf{Root-cause labels.} Every configuration that is not both fully Ready and smoke-test-passing receives one root-cause label from a fixed set of mechanisms (Table~\ref{tab:root-causes}). The label is \emph{observed} when the captured record shows the mechanism directly (for example, an error message naming it), \emph{inferred} when it is reconstructed from the manifest together with the observed status, and \emph{supported} when, in addition, configurations that differ from it only in the suspected setting succeed. Labels were assigned by one rater from the stored evidence; the label, mechanism and evidence for every configuration are in the replication package.

\textbf{D6 -- Recovery.} Two faults are injected into each configuration that serves, three times each, while an in-cluster prober calls the entry Service at 5~Hz: \emph{pod-delete} force-removes a Ready pod of the entry workload, and \emph{container-crash} stops its container at the node runtime. Metrics are recovery success, recovery time (from the fault to the replacement pod being Ready), and downtime (failed probes divided by the probing rate). Because a workload with one replica must be unavailable until its pod is replaced, while one with two or more need not be, D6 is reported by replica count.

No composite score is computed across dimensions.

\subsection{Evaluation pipeline and environment}\label{sec:method-pipeline}

Each configuration passes through the pipeline in a fixed order: parsing, schema validation, the D2/D3 scanners, the D4 checks, the server-side dry-run, deployment and smoke test (D5), and fault injection (D6). Failures are recorded as data: a parse failure marks every later stage as skipped; a schema failure does not block the static scanners or D4, which read the manifest text. D6 runs when the configuration is fully Ready or, failing that, when its smoke test passes, because an application can serve correctly while an auxiliary workload (for example, a load generator) is not Ready; D6 is also reported for fully Ready configurations only.

Every configuration is deployed into its own namespace, and the cluster is recreated between projects. Every image is served from a local registry mirror populated before deployment, so network latency does not confound time to Ready. Two deployment-time adjustments supply facts that the bundles do not contain, and both are applied identically to every source, recorded per configuration and never needed by a human manifest. First, for the projects whose exact image path appears only in unbundled build scripts, a reference that cannot be pulled as written but whose repository name identifies a documented component is served that component's canonical image; a reference that identifies no documented component fails. Second, a start-up prerequisite that an image enforces but the bundle never states is supplied when the configuration does not provide it: two environment variables in Online Boutique (one its front end requires, one that disables a cloud profiler in two services), and an empty Secret with the fixed name that the Kubernetes Dashboard expects. Permissions are never added. Section~\ref{sec:threats} bounds the effect of both adjustments.

The cluster has one control-plane node and three worker nodes, so that configurations using pod anti-affinity or a topology spread across three replicas can schedule. One setting departs from \texttt{kind}'s defaults: the container runtime limits each container to 1,048,576 open files (Docker's default), because some older images size internal tables from that limit and could otherwise exhaust the host's memory. Image tags are resolved when the mirror is populated, at evaluation time. A \texttt{LoadBalancer} Service is treated as valid but pending, an Ingress is created but not exercised, and PersistentVolumeClaims use \texttt{kind}'s local-path StorageClass, so a configuration that requests a StorageClass absent from this environment fails with that reason recorded.

\subsection{Statistical analysis}\label{sec:method-analysis}

Configurations are clustered: the three samples of a (model, project, prompt) cell share their input, and the projects differ widely. Inference therefore treats the project or the (model, project) cell as the unit wherever the question allows, and configuration-level tests are reported as descriptive support.

\emph{RQ1} compares each model's neutral-prompt configurations with the human manifests at configuration level, with the Kruskal--Wallis test \citep{kruskal1952ranks} followed by Dunn's test \citep{dunn1964multiple} for continuous metrics and Fisher's exact test for binary ones, and with Cliff's $\delta$ \citep{cliff1993dominance} as effect size (negligible $<0.147$, small $<0.33$, medium $<0.474$, large otherwise \citep{romano2006appropriate}). Because only seven human configurations exist, RQ1 is repeated at project level: each model's mean over the three samples of a project is paired with that project's human manifest. With $n\le7$ pairs a Wilcoxon test adjusted over five models cannot reach $p<0.078$, so this analysis reports Hodges--Lehmann differences with exact 95\% intervals. \emph{RQ3} compares the five models with a Friedman test \citep{friedman1937ranks} blocked on project, over the five projects every model completed, with Kendall's $W$ as effect size (with five blocks the $\chi^2$ approximation is coarse); configuration-level Kruskal--Wallis tests are reported alongside. Consistency is the across-sample standard deviation of each metric within a cell and the share of cells whose three samples agree on the runtime outcome. \emph{RQ5} pairs the two prompts at the (model, project) cell, comparing the cell's mean (or pass rate) under each prompt with the Wilcoxon signed-rank test; pairing individual samples is not meaningful, because the sample index carries no information across prompts. Exact McNemar tests \citep{mcnemar1947note} on configuration pairs are reported as a secondary view, per-model effects as Hodges--Lehmann estimates, heterogeneity between models with a Kruskal--Wallis test on the per-cell changes, and robustness by leaving out one project at a time. \emph{RQ4} models each binary runtime outcome with a logistic regression on one standardised static score, with prompt and project as fixed effects and standard errors clustered by (model, project) cell \citep{cameron2015cluster}; a second specification adds model fixed effects, and a cluster bootstrap over cells (999 replicates) checks the sandwich intervals. RQ2 reports frequencies, and agreement between Kubescape and Trivy on their overlapping checks is measured with Cohen's $\kappa$ \citep{cohen1960kappa}. Binary outcomes are shown with 95\% Wilson intervals \citep{wilson1927probable}.

$P$-values are adjusted with the Holm procedure \citep{holm1979simple} within explicitly defined families: the pairwise comparisons following one omnibus test; the five per-model tests of one metric; the three static scores (and, separately, the two deployment outcomes) of one prompt in RQ3; and the ten score--outcome models of RQ4.
